\documentclass[exercice,a4]{thedocument}%
\usepackage{pstricks}%
\usepackage{pst-plot}%
\usepackage{array}%
\startdatakeys
    \source{25GENMATAS1}
    \BO{2026}
    \cycle{4}
    \competencesDuSocle{CA,RA}
    \objectifsApprentissage{C4A1aa,C4A2ad,C4B2cb,C4D2aa,C4D3ac}
    \motsCles{facteurs premiers,trigonométrie,probabilité,fonction linéaire,image,développement}
\enddatakeys
\begin{document}%
\providecommand{\degres}{^\circ}%
    

Les 5 situations suivantes sont indépendantes. On rappelle que, sauf indications contraires, les réponses doivent être justifiées.

\subsection*{Situation 1}
Décomposer 390 en produit de facteurs premiers.

\subsection*{Situation 2}
ABC est un triangle rectangle en A avec AB = 10\,cm et BC = 20\,cm.

Quelle est la mesure de l'angle $\widehat{\text{ABC}}$ ?

\subsection*{Situation 3}
Une urne contient 12 jetons numérotés de 1 à 12 indiscernables au toucher. On pioche un jeton au hasard dans cette urne.

Quelle est la probabilité d'obtenir un nombre inférieur ou égal à 5 ?

\subsection*{Situation 4}
On considère la fonction $f$ dont on donne un tableau de valeurs et la représentation graphique ci-dessous :

\parbox{.4\textwidth}{\begin{center}
\begin{tabular}{|c|c|c|c|}
\hline
$x$ & 0 & 1 & 2 \\
\hline
$f(x)$ & 1 & $-1$ & $-3$ \\
\hline
\end{tabular}
\end{center}}\hfill\parbox{.5\textwidth}{\psset{unit=.5cm,arrowsize=2pt 3}
\begin{pspicture}(-3.7,-3.7)(3,8)
\psline[linewidth=2pt] (-3,7)(2.25,-3.5)
\psgrid[gridlabels=0pt,subgriddiv=0,gridwidth=0.15pt](-4,-4)(3,7)
\uput[ur](1,0){\scriptsize 1}
\psaxes[linewidth=1.25pt,labelFontSize=\scriptstyle](0,0)(0,1)
\psaxes[linewidth=1.25pt,labelFontSize=\scriptstyle,Dx=10,Dy=10]{->}(0,0)(-4,-4)(3,7)
\rput(0,8){\small Représentation graphique de la fonction $f$}
\end{pspicture}}

\medskip

\begin{enumerate}
\item Quelle est l'image de 2 par la fonction $f$ ? (\textbf{sans justifier})
\item Quelle est l'image de $- 1$ par la fonction $f$ ? (\textbf{sans justifier})
\item La fonction $f$ est-elle une fonction linéaire ?
\end{enumerate}

\subsection*{Situation 5}

On considère l'égalité suivante : $(2x - 3)(4x + 5) = 8x^2 - 2x - 15$.

\begin{enumerate}
\item Montrer que cette égalité est vraie pour $x = 2$.
\item Cette égalité est-elle vraie quelle que soit la valeur de $x$ ?
\end{enumerate}
\showthesource
\end{document}
